% Document class
\documentclass[12pt,a4paper]{article}

% Required Packages
\usepackage{amsmath, url}
\usepackage[francais]{babel}
\usepackage{fontspec}
\usepackage{color}
\usepackage{color,xcolor,ucs,soul}
\usepackage{sectsty}
\usepackage{hyperref, tikz}
\usepackage{eso-pic, titlesec, fancyhdr, fancyvrb}
\usepackage{multicol}
\usepackage{fancyvrb}
\usepackage{abstract}
\usepackage{listings}
\usepackage{framed,graphicx,xcolor}

\definecolor{shadecolor}{rgb}{1,0,0}

\providecommand{\tightlist}{%
  \setlength{\itemsep}{0pt}\setlength{\parskip}{0pt}}


% Configuration colors
\definecolor{dcblue}{HTML}{3A619B}
\definecolor{dcyellow}{HTML}{F4DD19}
\definecolor{dcbluelight}{HTML}{709EA4}
\definecolor{dcbluealt}{HTML}{25415E}
\definecolor{dcbluecover}{HTML}{335F93}
\definecolor{dcbluetext}{HTML}{002C5C}

% Font properties
\setmainfont{Gotham}
\defaultfontfeatures{Ligatures=TeX}
\AtBeginDocument{\color{dcbluetext}}
\partfont{\color{dcbluelight}}
\sectionfont{\color{dcbluelight}}
\subsectionfont{\color{dcbluelight}}
\subsubsectionfont{\color{dcbluelight}}

% Default commands

\newcommand{\twocols}[1]{
  \begin{multicols}{2}
    #1
  \end{multicols}
}

\fancyhead[R]{\bfseries\color{dcbluealt}\nouppercase\leftmark}
\fancyhead[L]{}
\renewcommand{\headrulewidth}{1.5pt}
\renewcommand{\headrule}{{\color{dcyellow}
\hrule width\headwidth height\headrulewidth \vskip-\headrulewidth}}
\fancyfoot[C]{\bfseries\color{dcbluelight}\thepage}

\newcommand{\VerbBar}{|}
\newcommand{\VERB}{\Verb[commandchars=\\\{\}]}
\DefineVerbatimEnvironment{Highlighting}{Verbatim}{commandchars=\\\{\}}
% Add ',fontsize=\small' for more characters per line
\newenvironment{Shaded}{}{}
\newcommand{\KeywordTok}[1]{\textcolor[rgb]{0.00,0.44,0.13}{\textbf{{#1}}}}
\newcommand{\DataTypeTok}[1]{\textcolor[rgb]{0.56,0.13,0.00}{{#1}}}
\newcommand{\DecValTok}[1]{\textcolor[rgb]{0.25,0.63,0.44}{{#1}}}
\newcommand{\BaseNTok}[1]{\textcolor[rgb]{0.25,0.63,0.44}{{#1}}}
\newcommand{\FloatTok}[1]{\textcolor[rgb]{0.25,0.63,0.44}{{#1}}}
\newcommand{\ConstantTok}[1]{\textcolor[rgb]{0.53,0.00,0.00}{{#1}}}
\newcommand{\CharTok}[1]{\textcolor[rgb]{0.25,0.44,0.63}{{#1}}}
\newcommand{\SpecialCharTok}[1]{\textcolor[rgb]{0.25,0.44,0.63}{{#1}}}
\newcommand{\StringTok}[1]{\textcolor[rgb]{0.25,0.44,0.63}{{#1}}}
\newcommand{\VerbatimStringTok}[1]{\textcolor[rgb]{0.25,0.44,0.63}{{#1}}}
\newcommand{\SpecialStringTok}[1]{\textcolor[rgb]{0.73,0.40,0.53}{{#1}}}
\newcommand{\ImportTok}[1]{{#1}}
\newcommand{\CommentTok}[1]{\textcolor[rgb]{0.38,0.63,0.69}{\textit{{#1}}}}
\newcommand{\DocumentationTok}[1]{\textcolor[rgb]{0.73,0.13,0.13}{\textit{{#1}}}}
\newcommand{\AnnotationTok}[1]{\textcolor[rgb]{0.38,0.63,0.69}{\textbf{\textit{{#1}}}}}
\newcommand{\CommentVarTok}[1]{\textcolor[rgb]{0.38,0.63,0.69}{\textbf{\textit{{#1}}}}}
\newcommand{\OtherTok}[1]{\textcolor[rgb]{0.00,0.44,0.13}{{#1}}}
\newcommand{\FunctionTok}[1]{\textcolor[rgb]{0.02,0.16,0.49}{{#1}}}
\newcommand{\VariableTok}[1]{\textcolor[rgb]{0.10,0.09,0.49}{{#1}}}
\newcommand{\ControlFlowTok}[1]{\textcolor[rgb]{0.00,0.44,0.13}{\textbf{{#1}}}}
\newcommand{\OperatorTok}[1]{\textcolor[rgb]{0.40,0.40,0.40}{{#1}}}
\newcommand{\BuiltInTok}[1]{{#1}}
\newcommand{\ExtensionTok}[1]{{#1}}
\newcommand{\PreprocessorTok}[1]{\textcolor[rgb]{0.74,0.48,0.00}{{#1}}}
\newcommand{\AttributeTok}[1]{\textcolor[rgb]{0.49,0.56,0.16}{{#1}}}
\newcommand{\RegionMarkerTok}[1]{{#1}}
\newcommand{\InformationTok}[1]{\textcolor[rgb]{0.38,0.63,0.69}{\textbf{\textit{{#1}}}}}
\newcommand{\WarningTok}[1]{\textcolor[rgb]{0.38,0.63,0.69}{\textbf{\textit{{#1}}}}}
\newcommand{\AlertTok}[1]{\textcolor[rgb]{1.00,0.00,0.00}{\textbf{{#1}}}}
\newcommand{\ErrorTok}[1]{\textcolor[rgb]{1.00,0.00,0.00}{\textbf{{#1}}}}
\newcommand{\NormalTok}[1]{{#1}}
\usepackage{graphicx,grffile}


% Document meta
\title{Rapport de stage TN09: Développement d'applications web en React et Redux}
\author{Aurélie \bsc{Digeon}}

% Document
\begin{document}

  {
    \begin{titlepage}
      \makeatletter
      \begin{tikzpicture}[
          overlay,
          remember picture,
          anchor=north west,
          inner sep=0pt,
          outer sep=0pt,
          fill opacity=.74
        ]
        \node[inner sep=0] at (current page.north west)
        {\includegraphics[height=\pdfpageheight]{layout/cover.jpg}};
        \draw[fill=dcbluecover]
          (current page.north west) rectangle (current page.south east);
      \end{tikzpicture}

      \begin{center}
        {\fontspec{Omnes}
          {\Huge {\textcolor{dcyellow}{\textbf{
            {\MakeUppercase{\@title}}
          }}\par}}
        }
        \vskip 2em
        {\large
        \begin{tabular}[t]{c}
          {\textcolor{white}{\@author}}
        \end{tabular}
        \par}
      \end{center}

      \vfill

      \begin{center}
        {\includegraphics[width=5cm]{layout/logo.png}}
      \end{center}
      \begin{center}
        {\includegraphics[width=3cm]{layout/logosUTC_SUnb.png}}
      \end{center}
      \makeatother
    \end{titlepage}
    \thispagestyle{empty}
    \tableofcontents
    \newpage
  }


  \pagestyle{fancy}
  %% start writing here

  \newpage

  \section{Remerciements}\label{remerciements}

  \bigskip

  Je tiens tout d'abord à remercier Benjamin Tierny, Robin Komiwes et
  Julien Vanden Torren de m'avoir acceuilli chez Dernier cri pour mon
  stage.

  \bigskip

  Je remercie également mon suiveur Harry Claisse pour son aide et son
  accompagnement, ainsi que les enseignants de l'Université technologique
  de Compiègne.

  \bigskip

  Merci également à toute l'équipe de Dernier cri pour avoir rendu mon
  stage si enrichissant et agréable et pour m'avoir accueuilli à bras
  ouverts.

  \bigskip

  J'aimerai finalement remercier mes parents et mon frère, qui m'ont
  permis de faire ses études et m'ont soutenu durant ce stage.

  \newpage

  \iffalse
  \# Résumé et abstract

  \bigskip

  \subsection{Résumé}\label{ruxe9sumuxe9}

  \newpage

  \subsection{Abstract}\label{abstract}

  \newpage

  \fi

  \section{Introduction}\label{introduction}

  \bigskip

  Dans le cadre de mon stage de TN09 lors de ma quatrième année
  d'ingénieur à l'Université technologique de Compiègne (UTC), j'ai
  effectué un stage de six mois chez Dernier cri.

  \bigskip

  Dernier Cri est une Start-Up créé en 2011 spécialisée dans l'innovation
  numérique. L'équipe est en charge du développement, du déploiement et de
  la maintenance d'applications pour le compte de plusieurs clients.

  \bigskip

  Ma mission a été d'intégrer l'équipe de développement pour aider dans la
  création de plusieurs applications web.

  \bigskip

  J'ai pu lors de ce stage intégrer une équipe dynamique et pro-active.
  J'ai notamment pu prendre part à de nombreuses présentations internes
  sur différentes technologies, à l'écriture d'article de blog. J'ai
  également pu participer activement à la relation client lors de mes
  projets.

  \bigskip

  Dans ce rapport je vais vous présenter tout d'abord Dernier cri,
  l'entreprise qui m'a acceuilli. Il me tient à coeur de vous exposer les
  techniques de travail et les différents moyens de s'exprimer qu'elle
  offre à ses équipes.

  \bigskip

  Je vous exposerai ensuite ma mission au sein de l'entreprise. Je
  détaillerai les technologies que j'ai utilisé mais aussi la gestion de
  projet ainsi que la relation client, qui sont des éléments clés dans la
  bonne conduite d'un projet.

  \bigskip

  Finalement, il me tient à coeur de présenter la communauté de Lille et
  ses particularités. En effet, le dynamisme de la communauté et la
  diversité des événements proposés m'ont beaucoup aidé à m'intégrer et à
  approfondir mon projet professionnel.

  \newpage

  \section{Dernier cri}\label{dernier-cri}

  \bigskip

  \subsection{Histoire}\label{histoire}

  \bigskip

  En 2011 Robin Komiwes et Benjamin Tierny créent Nectify dans le but de
  développer Fresc, un outil de partage d'avis sur des visuels. Bien que
  cet outil connait un succès certain avec aujourd'hui plus de 300
  sociétés utilisant Fresc à travers des milliers de projets, la
  rentabilité du projet n'est pas suffisante.

  \bigskip

  Nectify choisit alors de compléter ses revenus par de la prestation de
  services centrée sur l'innovation.

  \bigskip

  Début 2014, la majeure partie du chiffre d'affaire de Nectify était dû
  aux activités de prestations de services, Fresc ne représentant qu'une
  part marginale.

  \bigskip

  Devenant donc une agence spécialisée dans l'innovation digitale, Nectify
  choisit de créer sa propre image, distincte de Fresc. C'est dans ce
  mouvement que la société est devenue \textbf{Dernier cri}.

  \bigskip

  Aujourd'hui Dernier cri est une agence web qui met un point d'honneur à
  proposer à ses clients une solution complète adaptée à leur
  problèmatique spécifique. De la conception à la réalisation,
  l'entreprise accompagne ses clients de A à Z pour aboutir à un produit
  au plus proche des besoins de ceux-ci. Cela permet aux développeurs
  d'opérer dans différents domaines d'activités et d'avoir une vue globale
  du développement de produit.

  \bigskip

  \subsection{Secteur d'activité}\label{secteur-dactivituxe9}

  \bigskip

  Le secteur de l'informatique est aujourd'hui est pleine expension et ne
  connait pas la crise. Pour donner un ordre d'idée, le marché de la
  programmation et des services informatiques embauche près de 400 000
  personnes dans ses 21 000 ESN (Entreprises de services du numérique) et
  ne cesse de croître depuis ces cinq dernières années.

  \bigskip

  Dernier cri est une \emph{startup} spécialisée dans l'innovation
  digitale. Elle se démarque notamment des autres agence web en proposant
  plusieurs services, dont évidemment la création d'applications
  web/mobile, mais aussi l'application de la recherche fondamentale en
  apprentissage automatique et traitement de grosses données, afin
  d'augmenter la capacité d'innovation d'entreprises tiers. Cela se
  traduit notamment par de l'exploration de connaissances à partir de
  données pour des magasins, par exemple, ou la création d'audit.

  \bigskip

  C'est cette diversité qui permet à Dernier cri de se démarquer dans
  l'univers du web à Lille, où de nombreuse agence se partage le marché.

  \bigskip

  \subsection{Organisation}\label{organisation}

  \bigskip

  Dernier cri est une entreprise de petite taille et profite de cela pour
  mettre en place une organisation souple qui offre l'opportunité à tous
  de s'impliquer dans les projets, que ce soit des projets pour des
  clients ou des projets interne.

  \bigskip

  Au sommet de l'organisation de l'entreprise se trouve les deux
  fondateurs ainsi que le troisème associé. Le directeur général de
  l'entreprise (CTO) est Benjamin Tierny, tandis que Robin Komiwes est le
  directeur de la technologie (CTO). Le troisième associé, Julien Vanden
  Torren, est Account Manager, c'est à dire le lien entre l'entreprise et
  les clients.

  \bigskip

  L'entreprise compte aujourd'hui 18 employés mais se trouve dans une
  phase d'expension avec de nouvelles embauches en perspective, grave à de
  nombreux nouveaux projets.

  \bigskip

  L'entreprise est également constituée d'une chef de projet, Laetitia
  Cocusse, qui s'occupe de superviser la plupart des projets.

  \bigskip

  L'équipe de développeur est constituée notament d'un devops, c'est à
  dire d'une personne possédant à la fois les compétences d'un développeur
  et d'un ingénieur système (Jean-Serge Monbailly), d'un Data Scientist
  (Antonin Carette) qui travaille sur les projets de big data.

  \bigskip

  La plupart des développeurs travaille sur plusieurs projets en même
  temps, selon les besoins de l'entreprise et les compétences de chacun.
  Des équipes de 2/3 développeurs sont crées, mélangeant les compétences
  front et back, de développement et de système.

  \bigskip

  \subsubsection{Github et code review}\label{github-et-code-review}

  \bigskip

  L'entreprise utilise principalement la plateforme Github comme service
  web d'hébergement et donc le logiciel de gestion de versions Git. Github
  permet aux développeurs de travailler à plusieurs sur le même projet, de
  résoudre rapidement des conflits dus à la modification d'un même
  document par plusieurs personnes, de gérer différentes versions du
  projet et beaucoup d'autre problématiques.

  \bigskip

  Une nouvelle version de Github propose une section \textbf{Projet}
  permettant de gérer les \textbf{issues}, c'est à dire les tâches. Cette
  section permet notamment de séparer les tâches en plusieures colonnes,
  par exemple : \emph{A faire}, \emph{En cours}, \emph{Terminé}. Il est
  aussi possible d'attribuer les tâches à un contributeur, ou encore de
  leur attribuer des labels tel que \emph{Urgent}, \emph{Bug} ou encore
  une estimation de temps quand à la réalisation de la che.

  \bigskip

  Dernier cri utilisait jusque là un outil similaire. Il a été décidé
  d'utiliser la section \textbf{Projet} de Github pour les nouveaux
  projets, notamment ceux sur lesquels j'ai été affectée.

  \bigskip

  Dernier cri utilise ces outils mis à disposition par Github pour mettre
  en place un processus de vérification de la qualités du code et
  d'entraide. Chaque developpeur, une fois une tâche terminée, propose une
  \emph{Pull request}, c'est à dire demande à fusionner sa version du
  projet, modifiée pour résoudre la tâche, avec la version principale,
  stable. Il demande ensuite à ses collégues ayant des compétences dans le
  langage utilisé de relire et de commenter cette \emph{Pull request}.

  \bigskip

  C'est l'occasion pour les développeurs d'avoir l'avis de leurs collégues
  sur leur style d'écriture et leur façon de coder, ce qui permet souvent
  de découvrir de nouvelles méthodes et d'argumenter sur les meilleurs
  techniques à utiliser. Dernier cri utilise ce système pour garantir une
  certaine qualité du code ainsi qu'un style d'écriture homogéne.

  \bigskip

  \subsubsection{Talk interne}\label{talk-interne}

  \bigskip

  Dans l'optique d'un partage du savoir dans l'entreprise, les employés
  sont invités à faire des \emph{talks}, c'est à dire de petites
  présentations. Celles-ci ont pour objectif premier de présenter certains
  enjeux et solutions techniques en rapport avec la réalisation d'un
  projet, le résultat d'une veille, ou tout simplement un sujet qui les
  intéresse. Les \emph{talks} techniques permettent aux membres d'une
  équipe de partager entre eux leurs expériences et leur passions.

  \bigskip

  Ces présentations ont deux principaux avantages. Tout d'abord elles
  permettent aux spectateurs de la présentation d'apprendre de nouvelles
  choses. C'est l'occasion de découvrir un sujet dont on ne se doutait pas
  de l'intéret ou bien de l'existence. De plus, la présentation est
  travaillée et structurée, car une personne a déjà fait l'effort de trier
  les informations et de les présenter de la manière la plus lisible,
  claire, ou ludique. C'est ainsi souvent beaucoup plus facile d'aborder
  un sujet avec ces talks, que de partir à la pêche aux informations sur
  internet.

  \bigskip

  Le second avantage concerne l'organisateur de la présentation. Tout
  d'abord, sensibiliser ses coéquipiers à des problèmes peut être un
  véritable gain de temps pour le futur. En effet, apporter ses
  connaissances à l'équipe en entier permet d'éviter de devoir aller aider
  ou expliquer individuellement les même choses.

  \bigskip

  De plus, le rédacteur apprend beaucoup car, pour la création de la
  présentation, il doit devenir un expert sur le sujet qu'il va couvrir.
  Il doit appronfondir assez le sujet pour être capable, d'une part d'être
  clair dans sa présentation, d'autre part de répondrex aux éventuelles
  questions.

  \bigskip

  Pour lui c'est aussi une occasion de travailler sur ces compétences
  d'orateur. Cela peut évidemment lui être utile pour son travail, sa vie
  de tout les jours, mais surtout un techtalk interne peut faire office
  d'incubation pour une présentation dans un autre context. En effet, il
  existe à Lille de nombreux évènements où il est possible de faire des
  présentations, ce qui est très intéréssant mais aussi peut être assez
  intimidant. Le présentateur peut tester son talk en interne avant d'en
  faire partager la communauté.

  \bigskip

  Finalement, présenter un sujet issu d'une veille technologique permet,
  au-delà de faire découvrir aux autres un sujet potentiellement
  intéressant, de certifier l'intérêt du sujet. En effet, beaucoup de
  sujets qui au premier coup d'oeil semblaient prometteur, peuvent avec un
  peu de recherche se révéler creux et sans interet.

  \bigskip

  Les talks internes sont une très bonne manière de partager des
  connaissances de manière ludique. Durant mon stage j'ai eu l'occasion
  d'assister à des présentations sur de nombreuses technologies telles que
  React, Docker, Rust, MJML\ldots{} Mais aussi sur des sujets variés comme
  le Growth hacking (Ensemble de techniques de marketing permettant
  d'accélérer rapidement et significativement la croissance d'une
  start-up.), ou encore la fabrication d'un audit.. Ces présentations sont
  publiées sur la
  \href{https://www.youtube.com/channel/UCDfdBlzldhg_PEu3xZTPsHg}{chaîne
  youtube Dernier cri}. Cela m'a permis de découvrir de nouveaux sujets et
  de mieux comprendre les discutions, et les problématiques que
  rencontraient mes collégues.

  \bigskip

  \subsubsection{Article de blog}\label{article-de-blog}

  \bigskip

  Dernier cri posséde également un
  \href{http://derniercri.io/tech-blog}{blog technique} alimenté par les
  développeurs de l'équipe. Les objectifs sont notablement les mêmes que
  pour les présentations internes. Mais cela permet aussi à Dernier cri de
  rayonner et montrer ses compétences techniques et son esprit d'analyse.

  \bigskip

  Le blog technique permet à Dernier cri de se positionner en tant
  qu'expert. Publier des articles blog est excellent pour la réputation.
  En délivrant des conseils professionnels utiles, il prouve leur maîtrise
  de leur activité.

  \bigskip

  Cela permet également d'améliorer la visibilité de l'entreprise en
  exposant son savoir et savoir-faire à une large communauté
  professionnelle. Un blog est également utile pour tenir informé sa
  communauté : il permet de rester en contact permanent avec ses prospects
  et de les tenir informé de toute l'actualité de son entreprise, comme un
  nouveau produit, une nouvelle technologie maitrisée \ldots{}

  \bigskip

  Pour le rédacteur de l'article, cela a également de nombreux avantages.
  Par exemple, cela lui permet de travailler sa rédaction et son
  argumentation, de partager avec d'autres développeurs à travers les
  commentaires, de gagner en visibilité et en réputation, ou encore de se
  forcer à rester à la page, pour être exhaustif et à jour dans ses
  références.

  \bigskip

  J'ai moi-même eu l'occasion lors de mon stage d'écrire un article pour
  le blog technique de Dernier cri. Dans cette article je décris les
  technologies que j'ai utilisé lors de mon stage, React et Reduc, ainsi
  que cinq outils que j'ai découvert et qui m'ont facilité la vie pour
  utiliser ses technologies. Ce fut une expérience très enrichissante.
  Cela m'a poussé à enrichir mes connaissances dans le domaine, de
  clarifier les concepts pour savoir les expliquer à tous et de travailler
  ma qualité de rédaction.

  \newpage

  \section{Cadre du stage}\label{cadre-du-stage}

  \bigskip

  Durant mon stage j'ai rejoint l'équipe de développeurs de l'entreprise
  et j'ai pus participer au développement de deux applications. Ces deux
  projets s'appuyé sur du React, une bibliothèque JavaScript libre
  développée par Facebook depuis 2013. N'ayant jamais utilisé cette
  bibliothèque, j'ai donc du tout d'abord me former.

  \bigskip

  Le premier projet auquel j'ai participé se nomme Photolix. C'est un site
  internet de développement de photos, avec pour objectifs de toucher un
  large public et de limiter au maximum le temps d'attente du client en
  envoyant les photos au serveur dès leur sélection.

  \bigskip

  Le second projet est FinFrog, un site proposant des prêts financés par
  des particuliers. Ce projet était déjà assez avancé à mon arrivé. Le
  client possédait un site en ligne, mais souhaité changer l'apparence et
  ajouter des fonctionnalités, ce pourquoi il a fait appel à Dernier cri.

  \bigskip

  \subsection{Formation}\label{formation}

  \bigskip

  Lors de mon arrivée chez Dernier cri, j'ai eu l'occasion de me former
  sur Javascript ES6, React ainsi que Redux, car l'entreprise prévoyait de
  me mettre sur des projets utilisant ces technologie.

  \bigskip

  La popularité de ces trois technologie est en forte hausse, et de plus
  en plus demandé et utilisé pour la création d'application web. C'est
  pourquoi ce fut une véritable chance d'apprendre ces langages lors de
  mon stage.

  \bigskip

  Ma formation c'est faite à partir du site
  \href{https://www.codeschool.com/}{Code school}, disposant de cours en
  vidéos ainsi que d'exercices interactifs. Par la suite j'ai également
  pus compter sur Fabien Gavory, un développeur de Dernier cri, pour
  m'aider et m'expliquer certain concepts difficiles.

  \bigskip

  \subsubsection{Javascript ES6}\label{javascript-es6}

  \bigskip

  Le Javascript est un langage de programmation de scripts incontournable
  du web. Si à sa création il servait principalement à la réalisation
  d'animation, il est aujourd'hui au centre des applications. Le javacript
  sert maintenant à controler presque la totalité de l'application web.
  Cependant, ce langage n'étant pas prévu pour une telle complexité, il en
  résulte une syntaxe complexe et lourde. C'est dans ce contexte qu'un
  mise à jour du langage c'est imposé.

  \bigskip

  ES6 (ECMAScript Edition 6 ou encore ES2015) a été publiée en juin 2015.
  Il ajoute un ensemble de normes à celles déjà présente, pour apporter de
  nouvelle fonctionnalités qui permettent d'alléger le code, de le
  structurer, et de le rendre notamment plus maintenable, tout en restant
  compatibles avec le code existant.

  \bigskip

  ES6 n'est pas encore totalement supporté par les navigateurs, il est
  donc utilise d'utiliser un transcompilateur vers ES5, comme Babel.js.

  \bigskip

  L'apprentissage de ES6 a était primordiale pour mon stage : les
  nouvelles normes rende vraiment le code plus facile à lire et à écrire.
  J'ai ainsi dès le début de mon stage, pus prendre de bonne habitude
  quant au styles de mon code.

  \bigskip

  Cela m'a également permit de me familiariser avec les normes ECMAScript
  et de me persuader de la nécessité de rester attentive aux différentes
  actualités et évolution des langages. En effet, dans le milieu de
  l'informatique, les normes et bonnes pratiques sont très changeante et
  il est donc important de rester attentifs à l'actualité.

  \bigskip

  J'ai eu l'occasion durant mon stage de travailler sur un projet
  JavaScript n'utilisant pas ES6 et j'ai eu de grande difficultés à me
  passer des facilité d'écritures. Sans ES6, le code est beaucoup plus
  long et laborieux à lire et à écrire.

  \bigskip

  \subsubsection{React}\label{react}

  \bigskip

  Développé depuis 2013 par Facebook, \emph{React} est une bibliothèque
  JavaScript déclarative, efficace et flexible pour la création
  d'interfaces utilisateur. Cette bibliothèque c'est démarqué notamment
  par ses performances.

  \bigskip

  Elle est aujourd'hui utilisé par de nombreuses entreprises tel que
  Netflix, Yahoo, Airbnb ou encore Sony.

  \bigskip

  Une des particularités de \emph{React} est de découper l'application en
  composants, dépendant d'un état. Lors du changement de l'état d'un
  composant, \emph{React} génère les changements en \emph{HTML} pour les
  répercutés sur la page.

  \bigskip

  Cette bibliothèque est aujourd'hui en expansion. Elle a un succès
  certain auprès de la communautés des développeurs web, et de nombreux
  outils se développe autours. C'est donc un avantage d'avoir pu apprendre
  \emph{React} lors de mon stage, puis d'avoir mis en pratique ces
  connaissances lors des deux projets que j'ai effectués.

  \bigskip

  \subsubsection{Redux}\label{redux}

  \bigskip

  React, même s'il n'impose pas de bibliothèque pour les donnés et la
  communication des composants, offre une approche nommée \emph{Flux} très
  intéressante.

  \bigskip

  En plus de \emph{React}, Facebook à fournit une architecture appelé
  \emph{Flux} pour la gestion des donnés et la communication des
  composants. Cette architecture promet un flot unilatéral des données
  pour que le développeur puisse facilement suivre le trajet de des
  données d'un événement et ses conséquences.

  \bigskip

  \emph{Redux} est une des implémentation de \emph{Flux} les plus
  populaire, créée en Mai 2015 par Dan Abramov. Bien que reprenant les
  concepts de \emph{Flux}, \emph{Redux} les simplifie en utilisant des
  concepts liés à la programmation fonctionnelle.

  \bigskip

  La mise en place de \emph{Redux}, ou de \emph{Flux} en général peut
  sembler dans un premier temps laborieuse, car elle impose de nombreux
  fichier, de nombreux composants différents et une façon de penser qui
  peut désarçonner au début. Mais une fois les bases mises en place, cette
  architecture a l'avantage d'offrir une bonne lisibilité quant aux
  changements d'état. Il est possible de se retrouver dans le code très
  rapidement et d'ajouter des composants sans difficultés puisque
  \emph{React} et \emph{Flux} sont conçus pour être modulable. Flux trouve
  donc sont utilité dans les grandes applications.

  \bigskip

  Bien que \emph{Redux} soit la seule implémentation de \emph{Flux} que
  j'ai eu l'occasion d'utiliser, je pense qu'il s'agit d'une des
  meilleurs. Elle permet d'appréhender certain concept de la programmation
  fonctionnel, est simple à utiliser et assez populaire pour offrir un
  support en cas de problème.

  \bigskip

  \subsection{Gestion de projet}\label{gestion-de-projet}

  \bigskip

  Le processus de développement chez Dernier cri se passe généralement
  comme suit. Tout d'abord notre chef de projet, Laetitia Cocusse, discute
  avec le client pour comprendre ses besoins et ses attentes. Selon les
  besoins, elle organise une réunion avec le(s) developpeur(s) et le
  client pour clarifier certain point ou pour discuter des différents
  approches du problèmes possible. C'est l'occasion pour le développeur de
  proposer des solutions, peut être un peut différentes de celles
  imaginées par le client, ou encore de dire son avis sur les prochaines
  tâches.

  \bigskip

  Ensuite les taches sont estimées par un développeur. Il estime le temps
  qu'il pense passer sur le problème, en prenant en compte l'exploration
  de l'existant, le développement en lui-même, les tests et les possible
  retours. Cette estimation devra ensuite être validée par le client.

  \bigskip

  Une fois la tache validée, une \emph{issue} est créé dans Github avec
  une description, l'estimation et le développeur à qui est attribué la
  tache. Cela permet à chacun d'avoir une vue à tout moment de l'évolution
  du projet, des tâches en cours ou terminées. Il est également possible
  de commenter chaque \emph{Issue} pour, par exemple, demander des
  précisions, ou faire remonter une erreur.

  \bigskip

  Selon les priorités, le développeur choisi ou non l'ordre des taches.
  Pour protéger le projet actuel, stable, des nouvelles modification tant
  que celle-ci ne sont pas tester, le développeur crée une \emph{branche}
  dans Github, c'est à dire une copie du projet qui évoluera
  indépendamment de la branche principale. C'est sur cette branche que
  seront faites les modifications destinées a implémenter la nouvelle
  fonctionnalité.

  \bigskip

  C'est ensuite le moment de passer au développement à proprement parler.
  Le développeur utilise un enviromment de developpement sur son propre
  ordinateur pour simuler l'environnement de production. Cela peut se
  traduire par la création d'un base de donnée, ou la connexion à une
  \emph{API} spéciale. Il faut faire attention à ce que les données
  utilisée lors de la phase de développement n'ai pas d'impact sur celles
  de production.

  \bigskip

  Quand le développeur estime avoir terminer la tâche, il crée une
  \emph{Pull request} sur Github (comme expliqué plus haut dans le
  rapport). C'est alors le moment de prendre en compte les remarques et
  conseils des autres membres de l'équipe, de corrigé éventuellement des
  éléments, pour s'assurer de la qualité du code avant de l'incorporer
  dans le projet.

  \bigskip

  Une fois le développement de la fonctionnalité validé, la nouvelle
  version du site est déployé en \emph{Staging}, c'est à dire une version
  en ligne du site, qui à le même environnement que la production, mais
  avec des fausses données. Ce site sert à tester les nouvelles
  fonctionnalités avant de les pousser sur la production.

  \bigskip

  Laetitia fait une \emph{recette}, c'est à dire vérifie que la version de
  \emph{staging} remplie bien le besoin exprimé par le client, et que la
  nouvelle fonctionnalité n'a pas cassé autre chose. En cas de problème,
  le développeur reviens sur la tache jusqu'à ce que tout soit réglé.

  \bigskip

  Une fois un lot de taches effectuées, il est décidé en accord avec le
  client de poussé les modifications sur la production. Il faut alors
  vérifier que la \emph{mise en prodcution} c'est bien passé : que le site
  fonctionne toujours et que les nouvelles fonctionnalités sont bien en
  place.

  \newpage

  \section{Mes réalisations}\label{mes-ruxe9alisations}

  \bigskip

  -\textgreater{} Introduction

  \subsection{Photolix}\label{photolix}

  \subsubsection{Présentation du projet}\label{pruxe9sentation-du-projet}

  \bigskip

  Dès mon arrivée dans l'entreprise j'ai été assigné à la réalisation
  d'une application de développement photo. Le client possède un studio de
  développement photo sur Lille, et souhaitait proposé à sa clientèle un
  site simple et efficace.

  \bigskip

  Le principal objectif de ce projet était de télécharger les photos vers
  le serveur au fur et à mesure de leur sélection, pour ainsi éviter le
  temps d'attente du client à la fin de la saisie de ses informations.

  \bigskip

  Lors de mon arrivée sur le projet, un développeur de Dernier cri avait
  déjà posé des bases. Les fonctions de recadrage et de compression de la
  photo était notamment déjà écrite.

  \bigskip

  \subsubsection{Objectifs}\label{objectifs}

  \bigskip

  Les objectifs du projet Photolix étaient les suivants :

  \begin{itemize}
  \tightlist
  \item
    mettre en place le \emph{design} fournit par le client, à partir des
    maquettes sur Zeplin;
  \item
    gérer l'envoi des photos au serveur;
  \item
    mettre en place la possibilité de modifier les photos (formats,
    orientation\ldots{});
  \item
    pages de saisie des informations du client (adresses, informations de
    paiement) et page de remerciement;
  \end{itemize}

  \bigskip

  \subsubsection{Outils utilisés}\label{outils-utilisuxe9s}

  \bigskip

  Lors de mon arrivée sur ce premier projet, j'ai du apprendre à utiliser
  certains outils, tant au niveau de la gestion de projet, que du
  développement en lui-même.

  \bigskip

  Tout d'abord, le projet utilise Github comme service web d'hébergement
  et de gestion de développement, et par conséquent le logiciel de gestion
  de versions Git. Bien qu'ayant déjà utilisé Git et Github lors de mon
  DUT informatique, de projet personnel ou bien de projet à l'UTC, je ne
  connaissais pas certaines fonctionnalités de Github utilisé pars
  l'entreprise, notamment le code review et l'onglet projet. (Voir plus
  haut)

  \bigskip

  Le projet déjà existant utilisé npm comme gestionaire de paquets. npm
  est le gestionnaire de paquets officiel de Node.js. , automatiquement
  installé par défaut depuis la version 0.6.3 de Node.js. npm fonctionne
  avec un terminal et gère les dépendances pour une application. Il permet
  également d'installer des applications Node.js disponibles sur le dépôt
  npm. Il offre également la possibilité de créé des scripts. C'est une
  option vraiment pratique car grave a cela on peut construire et lancer
  l'application en une commande.

  \bigskip

  Pour mon environnement de travail, j'ai aussi utilisé un Linter. Code
  linting est un type d'analyse statique qui est fréquemment utilisé pour
  trouver des modèles problématiques ou le code qui ne respecte pas
  certaines directives de style. Il existe des linters de code pour la
  plupart des langages de programmation, et les compilateurs incorporent
  parfois le linting dans le processus de compilation.

  J'ai personnellement utilisé ESLint, qui est un utilitaire JavaScript
  open-source et libre créé à l'origine par Nicholas C. Zakas en Juin
  2013. ESLint est écrit en utilisant Node.js pour fournir un
  environnement d'exécution rapide et une installation facile via npm. La
  principale raison pour laquelle ESLint a été créée était de permettre
  aux développeurs de créer leurs propres règles de filtrage.

  \bigskip

  J'ai également pu utiliser le \emph{chatops} de Dernier cri, un outil
  d'administration système via la conversation. Intégré au Slack de
  l'entreprise, il permet à tout le personnel d'obtenir des informations
  sur un serveur ou une application et d'effectuer des résolutions simples
  en cas de panne. Concraitement, j'ai principalement utilisé le chatops
  pour déployer mon application.

  \bigskip

  L'application était écrite en React avec l'utilisation de Redux. J'ai
  pus donc mettre en application les principes appris lors de ma première
  semaine.

  \bigskip

  Pour l'intégration du style du site, j'ai pus utiliser Zeplin. C'est une
  application de collaboration pour les designers et les intégrateurs. Il
  permet aux designers de télécharger leurs maquette fonctionnelle
  directement à partir de Sketch et les ajouter aux dossiers de projet
  dans Zeplin. Les annotations seront automatiquement ajoutées aux designs
  (tailles, couleurs, marges et même suggestions CSS pour certains
  éléments). Il est alors beaucoup plus simple d'intégrer les maquettes.

  \bigskip

  Finalement, pour l'intégration des maquettes, j'ai fais le choix
  d'utiliser SASS. Sass (Syntactically Awesome Stylesheets) est un langage
  de génération dynamique de feuilles de style. On peut le voir comme une
  extension de CSS3, ajoutant de nouvelles règles dans notre façon
  d'intégrer un web design. Les principaux ajouts sont : les variables,
  les mixins, l'héritage de sélection et différents options très utiles.

  \bigskip

  \subsubsection{Déroulement}\label{duxe9roulement}

  \bigskip

  \paragraph{Intégration des
  maquettes}\label{intuxe9gration-des-maquettes}

  \bigskip

  La première partie du projet consisté à intégrer les maquettes fournit
  par le client. Dans un premier temps, j'ai du redécouper l'application.
  En effet, lors du premier jet réalisé par mon collègue, le client était
  parti sur une application monopage. Mais à la réception des maquettes,
  l'application était devenu multi-page. Il a donc tout d'abord mettre en
  place un routeur pour permettre à l'utilisateur de naviguer entre les
  pages.

  \bigskip

  J'ai choisi, après quelques recherches, d'utiliser React Router.
  \href{https://github.com/ReactTraining/react-router}{React Router} est
  une bibliothèque de routage pour React. Il dispose d'une \emph{API}
  simple avec des fonctionnalités puissantes. Il garde l'interface
  utilisateur synchronisé avec l'\emph{URL}. React Router est très simple
  à utiliser. Il suffit de lister les différentes routes souhaitées,
  associées au composent correspondants.

  \bigskip

  Une fois ce découpage effectuée, j'ai mis en place SASS pour la gestion
  des feuilles de style. Pour facilité son utilisation, j'ai créé
  plusieurs fichiers avec les variables et fonctions (mixins) qui seront
  utilisé dans tout le projet. Le fichier variable contient notamment les
  code hexadecimal des couleurs de l'application, les tailles des polices
  d'écriture utilisé, etc. Utilisé des variables évite de devoir revenir
  sur tous les fichiers du projet si l'on décide de changer l'une de ses
  variables.

  \bigskip

  Une fois ces fichiers SASS principaux créé, j'ai simplement créé un
  fichier SASS pour chaque composant React de l'application, qui contiens
  donc tout le style de ce composant. J'ai du faire beaucoup de recherches
  pour être à l'aise avec le CSS (Cascading Style Sheets), c'est à dire le
  langage décrivant la présentation de l'application. En effet, jusque là
  je n'avais pas eu l'occasion d'appronfondir mes connaissances en
  intégration.

  \paragraph{Téléchargement des
  photos}\label{tuxe9luxe9chargement-des-photos}

  \bigskip

  \paragraph{Modification des photos}\label{modification-des-photos}

  \bigskip

  \paragraph{Pages informations des clients, paiement et
  remerciement}\label{pages-informations-des-clients-paiement-et-remerciement}

  \bigskip

  \paragraph{Fin du projet}\label{fin-du-projet}

  \bigskip

  Vers mi-octobre, j'ai été réaffectée à un autre projet, laissant la fin
  du projet Photolix à Fabien. Il a fini de mettre en place la gestion du
  téléchargement des photos vers le serveur, notamment après leur
  modification. Il a également revu la fonction de redimensionnement des
  photos, car celle utilisé au debut n'était pas assez performantes : si
  l'on monté a une centaine de photos chargé, le navigateur ne supporté
  pas la charge.

  \bigskip

  J'ai pu étudier les modifications apporté par Fabien, et apprendre des
  erreurs que j'ai pus commettre. Par exemple, je n'avais pas assez
  travailler la gestion des erreurs. Il a fallut que Fabien reprenne mon
  travail pour ajouter l'affichage des erreur, en cas par exemple d'erreur
  lors du téléchargement des photos. Ces enseignements m'ont permis de ne
  pas reproduire ces erreurs dans le second projet sur lequel j'ai été
  affecté.

  \bigskip

  Finalement, j'ai pus retravailler sur le projet en decembre. Après un
  premier rendu au client, celui-ci souhaité quelques corrections ainsi
  que l'ajout de quelques fonctionnalité. Il avait utilisé un échantillons
  de client pour tester l'application et avait relevait des améliorations
  possible à l'interface utilisateurs. J'ai donc pus aider Fabien à mettre
  en place ces modifications.

  Developpement des fonctionnalités -\textgreater{} Détravailler

  Reunion client -\textgreater{} Détailler -\textgreater{} Client ouvert
  aux suggestion

  Collaboration avec Fabien -\textgreater{} Créateur du projet et maitrise
  le langage -\textgreater{} Très bon professeur -\textgreater{} Apprendre
  a communiquer à l'écrit + Channel Dev sur slack +

  \bigskip

  \subsubsection{Conclusion}\label{conclusion}

  -\textgreater{} Mettre en avant les difficultés et comment les surmonter
  -\textgreater{} Une chance d'avoir travailler tout de suite sur un vrai
  projet -\textgreater{} Good d'avoir eu qqn qui s'y connaissait (Fabien)
  pour m'aider -\textgreater{} Aprendre la gestion de projet de Dernier
  cri -\textgreater{} Laetitia -\textgreater{} Organisation

  \bigskip

  \bigskip

  \subsection{Finfrog}\label{finfrog}

  \subsubsection{Présentation du
  projet}\label{pruxe9sentation-du-projet-1}

  \bigskip

  Avec l'arrivée d'un nouveau client et la nécéssité de fournir un
  developpeur React sur ce projet, j'ai quitté le projet Photolix pour
  rejoindre Finfrog.

  \bigskip

  Finfrog est un projet de prêt collaboratif, c'est à dire que le site
  propose des prêt financés par des particuliers. Les prêts proposé vont
  de 200 à 600 euros, à rembourser en 1 à 3 mois. Le but de ce site est
  d'ouvrir, en acceptant des prêts qui ne serait pas validé par une banque
  car ils sont trop faibles ou bien que la personne est au chômage.

  \bigskip

  Lors de mon arrivée sur le projet, un site été déjà en ligne, développé
  par le client. Le premier objectif était de mettre un place un nouveau
  design sur ce site, d'abord sur la page d'acceuil, et ensuite sur les
  formulaire de demande de prêt.

  \bigskip

  Par la suite, j'ai été amené à developper de nouvelles fonctionnalités
  pour Finfrog, comme la partie du site réservée à la gestion des prêt par
  l'administrateur, les espaces emprunteur et préteur, la génération de
  contrat.

  \subsubsection{Objectifs}\label{objectifs-1}

  \subsubsection{Nouveaux outils}\label{nouveaux-outils}

  -\textgreater{} NPM -\textgreater{} On a commencé sur Bitbucket donc il
  a fallut apprendre -\textgreater{} Manipulation de la BDD
  -\textgreater{} On touche à la BDD -\textgreater{} pm2 -\textgreater{}

  \subsubsection{Déroulement}\label{duxe9roulement-1}

  \begin{itemize}
  \item
    Première partie design -\textgreater{} Page d'acceuil : prendre en
    main le projet et ses particuliarité -\textgreater{} Accés restraint
    au debut du projet -\textgreater{} Tunnel emprunteur : on rentre plus
    dans le techniques, la manipulation de redux
  \item
    génération de contrat
  \item
    Grosse passe pour le responsive + cross compatibilité des browsers
  \item
    Communication avec le client ? -\textgreater{} Surtout semaine de noel
    sans Laetitia
  \end{itemize}

  Citer le fait que je n'ai presque pas eu de code review a cause de la
  confidentialité

  \subsubsection{Conclusion}\label{conclusion-1}

  \newpage

  \section{La vie Lilloise}\label{la-vie-lilloise}

  Durant mon stage à Lille j'ai pus me rendre compte que la ville
  possédait une communauté web très active.

  \bigskip

  La ville à notamment reçu le label `French Tech' fin 2014, pour
  récompenser son dynamiste dans le numérique et l'innovation. Ce label,
  en plus de récompenser les efforts de la ville, constitue le point
  d'entrée vers des dispositifs nationaux comme des programmes pour
  attirer les entrepreneurs étrangers qui veulent créer leur
  \emph{start-up} en France.

  \bigskip

  La région profite de la présence de grands groupes nationaux comme
  Orange, Capgemini, IBM France, CGI, CISCO\ldots{}

  \bigskip

  De plus Lille a mis en place un ensemble de structures favorisant
  l'accompagnement et la croissance des startups vers un marché mondial.
  La plus notable est évidemment Euratechnologie, le Pôle d'excellence
  économique dédié aux Technologies de l'Information et de la
  Communication (TIC) de la métropole lilloise. EuraTechnologies a été
  classé dans le top 10 des accélérateurs de \emph{startup} d'Europe par
  Fundacity, et le 1er en France. Euratechnologie posséde des espaces
  dédié à la recherche, la formation et l'entrepreneuriat, un incubateur
  et un accélérateur.

  \bigskip

  La région lilloise posséde d'autre espaces dédié à l'innovation, les
  \emph{startup} et l'entrepreneuriat : La Plaine Images à Tourcoing et
  Roubaix, Eurasanté à Lille, La Haute Borne à Villeneuve d'Ascq, La Serre
  Numérique à Valenciennes, Le Pôle Numérique Culturel Louvre Lens Vallée
  de Lens\ldots{}

  \bigskip

  La région à connu l'émergence de nombreuse entreprises, prometteuse ou
  déjà fructueuse : Big Ben, Ankama, OVH, Addictiz, Stereograph, Clic \&
  Walk, Giroptic, Mazeberry, Vekia, Sparkow, Mdoloris, A-volute, Critizr,
  Intent Technologies\ldots{}

  \bigskip

  Avoir eu la chance de faire mon stage dans cette région m'a permis de
  profiter de cette écosysteme riche et actif. J'ai pus participer a des
  conférences, des salons et des réunions qui m'ont beaucoup apporté, tant
  au niveau technique que social. Cela m'a permit de préciser mon projet
  professionnel, en m'immergeant dans la vie entreprenariale d'une ville.

  \bigskip

  \subsection{Take Off Conference}\label{take-off-conference}

  \bigskip

  Dernier cri m'a donné l'occasion durant mon stage d'assister à la
  \href{http://takeoffconf.com/2016}{\textbf{Take Off Conference}} les 20
  et 21 octobre 2016. Cet évènement à lieu depuis plusieurs années à
  EuraTechnologies.

  \bigskip

  Historiquement, ce sont les fondateurs de Dernier cri qui ont créé la
  \textbf{Take Off Conference}, avec Florian Le Goff. Aujourd'hui ce sont
  d'autres acteurs de la communauté web de Lille qui ont pris le relais
  pour proposer une nouvelle édition.

  \bigskip

  La \textbf{Take Off Conference} est un cycle de conférences anglophones.
  L'événement dure 2 jours, et accueille des conférenciers du monde
  entier. Bien qu'elle reste avant tout une conférence pour les
  développeurs Web, elle reste accessible pour les développeurs en
  général.

  \bigskip

  Dernier cri m'a permis de participer à cette conférence avec l'un de mes
  collègues. Ce fut une véritable chance pour moi de rencontrer et
  échanger avec des conférenciers du monde entier. Les conférences étaient
  très intéréssantes et inspirantes, sur des sujets très varié allant de
  la compréhension des enjeux de mise en place de nouveau outils, à des
  sujets plus sociaux comme l'acceuil des developpeurs anglais après le
  Brexit.

  \bigskip

  Des évènements ont également été organisé le soir pour permettre au
  participant et au conférenciers d'échanger dans un cadre puis détendue.

  \bigskip

  Ce fut une excellente occasion de découvrir de nouvelle technologie, de
  m'ouvrir à des problématique que je ne connaissais pas ainsi que de
  rencontrer des développeurs qualifié et passionné. Lors de ses échanges,
  j'ai pus me rendre compte de la portée internationales de la
  programmation : des personnes des quatres coins du monde se retrouver
  sur les mêmes problématiques.

  \bigskip

  \subsection{Meetup}\label{meetup}

  \bigskip

  La communauté web de Lille est très active pour organisé des évenements.
  De très nombreux Meetup sont organisé sur différent domaine du web,
  acceuillant autant des experts que des débutant ou simplement des
  curieux.

  \bigskip

  Très actif dans cette communauté, Dernier cri acceuille au sein de ses
  locaux certaines de ses rencontres. J'ai notamment eu l'occasion
  d'assister à des Meetup de Lille FP (Fonctionnal programming, c'est à
  dire programmation fonctionnelle) et de Lille Elixir.

  \bigskip

  Certain membre de l'équipe sont également investi dans ces rencontres,
  en tant qu'organisateur ou bien speaker. Les patrons les poussent à
  (prendre part a la vie de la communauté web). J'ai ainsi pu facilement
  être au courant des différent évenement Lillois et y participer avec mes
  collégues, ce qui m'a permis d'être bien intégré.

  \bigskip

  \subsection{Maker Faire}\label{maker-faire}

  \bigskip

  La Maker Faire est un autre événement majeur organisé à Lille durant mon
  stage. Ce concept totalement unique regroupe stands de démonstration,
  ateliers de découverte, spectacles et conférences autour des thèmes de
  la créativité, de la fabrication, des cultures Do It Yourself et Makers.

  \bigskip

  Cet événement, présenté par Leroy Merlin en partenariat avec la Ville de
  Lille et lille3000 (programme culturel de la ville de Lille), réunit des
  passionnés de technologies, des artisans, des industriels, des amateurs,
  des ingénieurs, des clubs de science, des artistes, des étudiants et des
  Start'Up. Ils forment la communauté des Makers et viennent pour montrer
  leurs créations, partager leurs connaissances\ldots{}

  \bigskip

  Dernier cri à été invité à assister à la Maker Faire, notamment car nous
  développons l'application web de TechShop, l'atelier collaboratif de
  Leroy Merlin. Dans ce contexte j'ai pu découvrir la communauté
  \emph{Maker} de Lille. J'ai pus notamment décrouvrir les \emph{repair
  coffee}, lieux où l'on peut amener ses appareil électronique cassé et
  recevoir de l'aide pour leur réparation. Il y avait également des
  robots, des imprimantes 3D, des casques de réalité augmenté \ldots{}
  C'était un endroit plein d'aspiration et d'envie d'entreprendre.

  \newpage

  \section{Conclusion}\label{conclusion-2}

  -\textgreater{} Compétences apprises : Node React Redux ES6 Relation
  client Git Github pm2 npm

  -\textgreater{} Important : relation client et autonomie, apprentissage
  rapide, equipe disponible et polivalente (qui peut t'aider sur tout)

  -\textgreater{} environnement cool : Ambiance, tech off conf , meetup,
  plus belle vue de Lille

  \newpage

  \section{Glossaire ?}\label{glossaire}


\end{document}
